\inicializarSubseccion{Analisis de la novena pregunta (Por parte de Andrés)}
\textbf{9. ¿Qué tan grande debería de ser el diámetro del portavasos?}
Similar al séptimo caso, es fácil ver que los datos son completamente independientes de sí mismos. Otra manera de verlo es hacer una diferencia entre las diferentes probabilidades que tenemos al escoger uno de los datos de manera aleatoria. Al haber una equivalencia entre la cantidad de elementos en el universo y ver que no hay uniones y/o intersecciones entre dichos elementos. Es fácil ver que se cumple el lema de eventos mutuamente independientes, ósea:

\begin{equation}
P(A_1\cap \dots \cap A_n)=P(A_1)\times \dots \times P(A_n) 
\end{equation}

Al tener seis elementos independientes, con este análisis podemos ver que la moda, con 36 elementos dentro del conjunto (representando 35.6\% de la cantidad de elementos dentro del universo) es el dominio 8.50-10.50 cm. Es por eso que el diámetro del portavasos debería de ser entre 8.50 a 12.50 cm, de esta forma podemos abarcar la mayoría de los casos posibles, y así satisfacer a la mayoría de los consumidores.